\documentclass[a4paper]{article}
\usepackage{graphicx}
\usepackage{amsmath,amssymb}
\usepackage{hyperref}
\usepackage{minted}
\usepackage{multirow}

\title{Methodology}
\date{}
\begin{document}
    \maketitle
    \section{Embodied Cognitive Architecture}
    
    Using one of the already stored ``memory>longterm>episodic>autobiographical>Event specific knowledges`` to learn the temporal dependency between different modalities of an ``event specific knowledge`` and then compare it to the event it is experiencing at present time to detect any source that raises suprise in the agent. There are three suprise sources.
    Anomaly, outlier and novelty. Eventhough the framework we will introduce is capable of detectiong all, but we insiste on novelty detection.
    
    If all what an agent knows about the environment is shown by
    [[
        P(t)
    ]]
\end{document}